\section{Planned Extensions}
\label{sec:roadmap}

This report is a living document, and this section records what is expected
to change. Items are listed as agreed priorities or as proposals under
consideration; the changelog in Appendix~\ref{app:changelog} records what has
actually been delivered in each version.

\subsection{Database}

\paragraph{Additional economies.}
Extending coverage to further economies --- house prices, personal income and
rents --- is the current priority for the database. Candidates under
discussion include Slovakia and China. Each addition requires a national
source that can be reconciled with the FHFA benchmark on the criteria in
Section~\ref{subsec:construction}, and historical series sufficient to extend
it back towards the start of the sample; the constraint is generally the
availability of the historical series rather than the current one.

\paragraph{Nominal series in local currency.}
All series are currently published as indices on a 2005 base, which supports
comparison of dynamics but not of levels, and which makes reported figures
harder to interpret in absolute terms. Publishing nominal price data in local
currency alongside the indices would make the figures more directly readable,
at the cost of introducing a series that is explicitly not comparable across
countries. The presentational separation between the two would need care.

\subsection{Econometric Toolkit}

\paragraph{Extending the nowcast.}
The nowcast model of Section~\ref{subsec:nowcast} is currently specified for
a limited set of economies. Extending it to the United Kingdom and Spain is
proposed as the next step. The binding requirement is monthly housing
indicators of comparable coverage and timeliness to those used in the
existing specification; where the available indicators differ, the
specification will need to be adapted rather than transferred directly.

\paragraph{Extending the fundamentals decomposition.}
The PSY-IVX decomposition covers fewer economies than the exuberance tests,
limited by the availability of rent and long-term interest rate series
(Section~\ref{subsec:psyivx}). Widening that coverage would allow the
fundamentals question to be asked for a larger part of the panel.

\paragraph{Additional models.}
Further additions to the toolkit have been raised in principle. These are at
the stage of establishing whether there is appetite for broadening the
Observatory's methodological scope, rather than of specific proposals, and
will be recorded here as they are agreed.

\subsection{Real-time Evaluation}

Both the nowcast validation (Section~\ref{subsec:nowcast}) and any assessment
of the exuberance indicators currently use revised data. Because IHPD
observations are revised retrospectively
(Section~\ref{subsec:limitations}), and because archived vintages are
available through \texttt{ihpd\_get()}, a genuine real-time evaluation is
feasible and would give a more honest measure of how the Observatory's
indicators perform at the moment they are published.
